\section{回帰 : 関係を直線で表し, 予測に使う}

\subsection{問い : 広告費を増やしたら, 売上はどのくらい変わるか}

前章では, 2つの変数の関係の方向と強さを相関係数で測った.
広告費と売上の相関係数は$r = 0.87$だった.
強い正の相関がある. 広告費が多い月は売上も多い傾向がある, というところまではわかった.

しかし, 本当に知りたいのはもう少し先のことだ.
広告費を10万円増やしたら, 売上はどのくらい変わるのか.
来月の広告費が決まったとき, 売上をいくらと見込めばよいのか.
相関係数はこの問いには答えてくれない.
$r = 0.87$は「強く連動している」とは教えてくれるが, 「いくつ変わるか」は教えてくれないからだ.

なぜ答えられないのか.
$r$は$-1$から$1$の範囲に収まる単位のない数値だ.
「$x$が1増えたときに$y$がどれだけ動くか」という情報を含んでいない.
相関係数に欠けているのは\textbf{傾き}の情報なのだ.

散布図の点の並びに直線を当てはめれば, その傾きから「$x$が1増えたとき$y$が平均的にいくつ変わるか」を読み取れる.
さらに, まだ見ていない$x$の値に対して$y$を予測することもできる.

直線を引けばよいという方針は見えた.
では, その直線はどうやって決めるのだろうか.


\subsection{素朴な方法と限界 : 目で引いた線は当てになるか}

散布図に直線を引くこと自体は, 定規さえあればできる.
右上がりに並んだ点を眺めて, 「だいたいこのあたりかな」と線を引く.
日常の場面では, 実際にこれで十分なこともあるだろう.

しかし, 同じ散布図を3人に渡して, それぞれ「最も良い」と思う直線を引いてもらったらどうなるか.
3人の線はまず一致しない.
ある人は急な傾きの線を引き, 別の人はゆるやかな線を引く.
右端に離れた点がある場合, そこを重視する人もいれば, 中央に密集した点に引きずられる人もいる.

傾きが1.5倍も違う線が出てくることは珍しくない.
一方の線からは「広告費10万円増で売上5万円増」と読め, もう一方からは「売上8万円増」と読める.
同じデータから60\%も違う結論が出てしまう.
自分ならどこに線を引くか想像してみると, 他の2人と一致する自信はあまりないはずだ.

なぜこうなるのか.
目分量で線を引くとき, 人は散布図のどの部分に注目するかを無意識に選んでいる.
右端の点に目が行く人, 中央の密集に合わせる人, 全体の「雰囲気」で決める人.
注目する場所が違えば, 出てくる線も違う.

つまり, 問題は2つある.
\textbf{基準がない}ことと, \textbf{再現性がない}ことだ.
「良い直線」とは何かの定義がなければ, 誰が引いても同じ線にはならない.
同じ人がもう一度引いたとしても, 微妙に違う線になるかもしれない.

目分量に頼らず, 誰が計算しても同じ1本の線が出る方法が必要だ.


\subsection{直線の当てはめと評価 : 最小二乗法, $R^2$, 残差}

\subsubsection*{直線の式 : $y = a + bx$}

2変数$x$と$y$の間に直線的な関係があると想定して,
\[
  y = a + bx
\]
という形でその関係を表すことを考える.

$b$は\textbf{傾き}だ.
$x$が1増えたとき, $y$が平均的にいくつ変わるかを表す.
広告費と売上の例なら, $b = 2.3$は「広告費1万円増で売上が平均2.3万円増える」と読める.

$a$は\textbf{切片}で, $x = 0$のときの$y$の値にあたる.
ただし, $x = 0$が現実的な意味を持たない場合も多い.
広告費がゼロの月が実際にはありえないのなら, 切片の数値そのものにこだわる必要はない.
直線の位置を決めるために数学的に必要な部品だと思えばよい.

問題は, この$a$と$b$をどう決めるかだ.

\subsubsection*{残差 : 直線からのずれを測る}

もしデータの点がすべて1本の直線の上に乗っていたら, $a$と$b$は簡単に決まる.
どの2点を使っても同じ直線が出るからだ.
しかし現実のデータは, そうはいかない.
どんな直線を引いても, 実際の値と直線が予測する値のあいだにはずれが残る.

このずれを数値として扱えるようにしたい.
$i$番目のデータ点$(x_i, y_i)$に対して, 直線が予測する値は$a + bx_i$だ.
実際の値$y_i$とこの予測値の差を\textbf{残差}と呼び, $e_i$と書く.

\begin{defi}{残差}{residual}
データ点$(x_i, y_i)$と回帰直線$y = a + bx$に対し,
\[
  e_i = y_i - (a + bx_i)
\]
を$i$番目の\textbf{残差}と呼ぶ.
\end{defi}

$e_i$が正なら, 実際の値が直線より上にある.
負なら, 直線より下にある.
すべての点を通る直線は引けない以上, 残差をゼロにすることはできない.

ならば, 残差をなるべく小さくする直線が「良い直線」だろう.
では, 「残差が全体として小さい」とはどういう状態だろうか.

残差をそのまま合計してみるとどうなるか.
正の残差と負の残差が打ち消し合って, 合計はゼロに近くなることがある.
合計がゼロだからといって, 直線がデータに合っているわけではない.
打ち消しが起きているだけだ.

では打ち消しを防ぐにはどうすればよいか.
第2章で分散を定義したとき, 同じ問題に出会ったのを覚えているだろうか.
偏差の合計がゼロになるのを防ぐために, 二乗してから足したのだった.
ここでも同じ手を使う.
各残差を\textbf{二乗}してから合計する.

\subsubsection*{最小二乗法 : 残差の二乗和を最小にする}

この「残差の二乗和を最小にする$a$と$b$を選ぶ」という方針に名前がついている.

\begin{defi}{最小二乗法}{ols}
$n$個のデータ$(x_1, y_1), \ldots, (x_n, y_n)$に対し,
\[
  S(a, b) = \sum_{i=1}^{n} e_i^2 = \sum_{i=1}^{n} (y_i - a - bx_i)^2
\]
を最小にする$a$と$b$を選ぶ方法を\textbf{最小二乗法}（OLS: Ordinary Least Squares）と呼ぶ.
\end{defi}

$S(a, b)$は残差の二乗和だ.
$a$と$b$の値が変われば$S$の値も変わる.
この$S$を最小にする$a$と$b$は, 計算によって一意に求まる.

\[
  b = \frac{\displaystyle\sum_{i=1}^{n}(x_i - \bar{x})(y_i - \bar{y})}{\displaystyle\sum_{i=1}^{n}(x_i - \bar{x})^2}, \qquad a = \bar{y} - b\bar{x}
\]

見慣れない式に見えるかもしれないが, 部品は前章で出会ったものばかりだ.
分子の$\sum(x_i - \bar{x})(y_i - \bar{y})$は, $x$と$y$の偏差の積の合計であり, 前章で学んだ共分散の分子と同じ構造をしている.
分母の$\sum(x_i - \bar{x})^2$は, $x$の偏差の二乗和で, $x$の分散の分子に対応する.
つまり傾き$b$は, $x$と$y$がどれだけ一緒に動くか（共変動）を, $x$自身の散らばり（変動）で割ったものだ.
「$x$が1単位散らばるごとに$y$が平均的にどれだけ動くか」という直感そのままの形になっている.

切片$a = \bar{y} - b\bar{x}$は, 回帰直線が必ず点$(\bar{x}, \bar{y})$を通ることを意味する.
データの重心を通り, 傾き$b$で引いた直線.
そう読むと, 直線の位置が自然に定まっていることがわかる.

ところで, なぜ絶対値ではなく二乗なのかと思うかもしれない.
絶対値を使う方法（最小絶対偏差法）も存在するし, 外れ値に対してはむしろそちらのほうが頑健だ.
しかし, 二乗のほうが数学的に扱いやすく, 解が上の公式のように閉じた形で求まるという利点がある.
第2章で分散を定義したときとまったく同じ判断だ.
「打ち消しを防ぐために二乗する」という道を選んだおかげで, 計算が公式に乗る.
ばらつきの測り方として人間が決めた約束であって, 自然にそうなるわけではない.
ただ, その約束が実用上きわめて便利なのだ.

\begin{mybox}{ガウスとケレスの軌道}
1801年, イタリアの天文学者ピアッツィが小惑星ケレスを発見した.
しかし観測できたのはわずか41日間で, ケレスは太陽の向こう側に隠れてしまった.
再び姿を現したときにどこにいるかを, 不完全な観測データから予測しなければならない.
ガウスは観測値と予測値のずれの二乗和を最小にする方法で軌道を推定し, その予測は的中した.
ケレスは予測された位置で再発見された.

最小二乗法は, 「手元にあるのは不完全なデータだが, そこから最も確からしい答えを引き出したい」という切実な動機から実用化された道具だ.
なお, フランスのルジャンドルが1805年に最小二乗法を先に公表しており, 優先権を巡る論争がある.
ガウスは「自分は1795年から使っていた」と主張したが, 出版が後だったため決着はついていない.
\end{mybox}


\subsubsection*{相関係数との関係}

傾き$b$の式と前章の相関係数$r$の式を見比べると, 次の関係が導かれる.
\[
  b = r \cdot \frac{s_y}{s_x}
\]

$r$は関係の強さと方向を$-1$から$1$で表す, 単位のない指標だった.
$s_y / s_x$は, $y$の散らばりが$x$の散らばりの何倍かを表している.
$r$に$s_y / s_x$を掛けると, 単位のない「連動の度合い」が, 「$x$が1単位増えたとき$y$がいくつ変わるか」という単位つきの量に変換される.

たとえば$r = 0.87$, $s_y = 30$万円, $s_x = 15$万円なら,
\[
  b = 0.87 \times \frac{30}{15} = 1.74 \text{（万円/万円）}
\]
となる.
「強い相関がある」という抽象的な評価が, 「広告費1万円あたり売上1.74万円」という具体的な変化量に変わった.
相関係数と回帰の傾きは, 同じ直線関係を異なる角度から見た量なのだ.

直線の引き方と, 相関係数との関係がわかった.
では, 引いた直線はデータをどのくらいうまく捉えているだろうか.

\subsubsection*{$R^2$ : 説明できた割合}

直線を引いた.
では, この直線はデータをどのくらいうまく捉えているのか.

もし$x$の情報がまったくなかったとしたら, $y$の値をどう予測するだろうか.
使えるのは$\bar{y}$, つまり$y$の平均しかない.
どのデータ点に対しても$\bar{y}$と答えるしかないから, 予測は当然はずれる.
各点と$\bar{y}$のずれの二乗和$\sum(y_i - \bar{y})^2$が, $y$の全体の散らばりだ.

回帰直線を引くと, 予測値が$\bar{y}$から$\hat{y}_i = a + bx_i$に変わる.
$x$の情報を使ったぶんだけ, 予測は改善されるはずだ.
では, どのくらい改善されたのか.
「全体の散らばりのうち, 直線で説明できた割合」を測るのが\textbf{決定係数}$R^2$だ.

\begin{defi}{決定係数}{r-squared}
$y$の全変動を$\mathrm{SS_{total}} = \sum(y_i - \bar{y})^2$, 残差による変動を$\mathrm{SS_{res}} = \sum(y_i - \hat{y}_i)^2$とする.
決定係数は
\[
  R^2 = 1 - \frac{\mathrm{SS_{res}}}{\mathrm{SS_{total}}}
\]
で定義される. $\hat{y}_i = a + bx_i$は回帰直線による予測値である.
\end{defi}

$\mathrm{SS_{res}} / \mathrm{SS_{total}}$は, 直線で説明しきれなかった散らばりの割合だ.
これを1から引いているから, $R^2$は「説明できた割合」になる.

$R^2 = 0.76$なら, $y$の散らばりの76\%が$x$との直線関係で説明でき, 残り24\%は直線では捉えきれていない.
$R^2 = 1$なら全点が直線上に乗っており, $R^2 = 0$なら$x$と$y$に線形関係がない.

単回帰では, $R^2$は相関係数$r$の二乗に一致する.
$r = 0.87$なら$R^2 = 0.87^2 \approx 0.76$だ.
前章で「相関が強い」と評価したことと, 本章で「76\%説明できた」と評価することが, 数値としてつながっている.

\subsubsection*{$R^2$だけでは安心できない}

$R^2$が高ければ直線がデータをよく捉えている, と言いたくなる.
しかし, $R^2$は全体の当てはまりを1つの数値に圧縮した指標だ.
圧縮する過程で, ずれ方のパターンが見えなくなる.

アンスコムのカルテットという有名な例がある.
4組のデータ（各11点）のすべてで, $r \approx 0.82$, $R^2 \approx 0.67$, 傾きも切片もほぼ同じだ.
数値だけ見たら, 4組とも「まあまあの直線関係」で片づけてしまうだろう.

ところが散布図を描くと, 実態はまるで違う.

% [図挿入: アンスコムのカルテットの4枚の散布図]

1組目は素直な直線関係で, 回帰直線がデータをよく捉えている.
2組目は完全な放物線で, 直線を当てはめること自体が不適切だ.
3組目はほぼ完璧な直線に外れ値が1点だけあり, その1点が傾きを歪めている.
4組目は$x$が1点だけ離れた位置にあり, その1点が見かけの相関を作り出している.

1973年にフランシス・アンスコムがこの4組を発表した意図は明確だ.
統計量を計算する前にグラフを描け, ということだ.
$R^2$や傾きの数値だけでは, データの本当の姿はわからない.
直線が妥当かどうかを判断するには, 残差の様子を自分の目で確認する必要がある.

\subsubsection*{残差プロット : 直線の妥当性を目で確かめる}

残差$e_i = y_i - \hat{y}_i$を, 予測値$\hat{y}_i$に対してプロットした図を\textbf{残差プロット}と呼ぶ.

直線がデータの関係をうまく捉えていれば, 残差はゼロのまわりにランダムに散らばる.
特定のパターンは見えないはずだ.
逆に, 残差に系統的なパターンが見えたら, 直線では捉えきれない構造がデータに残っている.

たとえば残差がU字型に並んでいれば, 本当は曲線的な関係なのに直線で無理に近似した可能性がある.
残差の散らばりが予測値に応じて広がっていくなら, $x$が大きい領域で予測が不安定になっている.

残差プロットで確認するのは2点だ.

\begin{itemize}
  \item \textbf{パターンがあるか}. U字型, 扇形, 周期的なうねりなどが見えたら, 直線モデルの前提を疑う.
  \item \textbf{飛び離れた点があるか}. 他の残差から大きく外れた点があれば, そのデータ点を個別に確認する.
\end{itemize}

パターンがなく, ゼロのまわりにランダムに散らばっていれば, 「この直線はデータの関係をおおむね捉えている」と判断してよい.
厳密な検定もあるが, 本書では目視で十分だ.

これで, 直線を引く道具（最小二乗法）と, 引いた直線を疑う道具（$R^2$と残差プロット）の両方が揃った.


\subsection{手法 : 回帰直線を求めて評価する}

\subsubsection*{前提条件}

回帰直線を当てはめる前に, 以下の前提を確認する.

\begin{itemize}
  \item 2変数の間に\textbf{直線的な関係}があること. 散布図で確認する.
  \item 各データ点が互いに\textbf{独立}であること. データの集め方に依存する.
  \item 極端な外れ値がないこと. あるいは, その影響を把握していること.
\end{itemize}

直線性は散布図を見ればよい.
明らかに曲線的な関係しか見えないなら, 直線を当てはめるのは不適切だ.
独立性はデータの集め方で決まるため, 散布図からは判断しにくい.
たとえば月次の売上データなら, 前月の売上が翌月に影響していないかを考えておく必要がある.

\subsubsection*{手順}

\textbf{Step 1. 散布図を描く.}
$x$と$y$の散布図を描き, 直線的な関係がありそうかを目で確認する.
極端な外れ値がないかもこの段階で見ておく.

\textbf{Step 2. 傾き$b$と切片$a$を計算する.}
\[
  b = \frac{\sum(x_i - \bar{x})(y_i - \bar{y})}{\sum(x_i - \bar{x})^2}, \qquad a = \bar{y} - b\bar{x}
\]

\textbf{Step 3. $R^2$を計算する.}
\[
  R^2 = 1 - \frac{\sum(y_i - \hat{y}_i)^2}{\sum(y_i - \bar{y})^2}
\]
$y$の散らばりのうち, 直線で説明できた割合を確認する.

\textbf{Step 4. 残差プロットを描いて確認する.}
残差$e_i = y_i - \hat{y}_i$を予測値$\hat{y}_i$に対してプロットし, 系統的なパターンがないかを確認する.
パターンがなければ, 直線モデルはおおむね妥当だと判断できる.

この手順を, 前章から引き継いだ広告費と売上のデータに当てはめてみる.


\subsection{実践 : 広告費と売上で手を動かす}

前章で$r = 0.87$を計算した広告費と売上のデータを, そのまま使う.
ある店舗の10ヶ月分のデータだ.

% TODO: 具体的な数値テーブルはランニング例確定後に挿入
% 以下は仮のデータで計算手順を示す

\subsubsection*{Step 1 : 散布図の確認}

% [表挿入: 月, 広告費(万円), 売上(万円) の10行]
% [図挿入: 広告費 vs 売上の散布図]

散布図を描くと, 右上がりの傾向が見える.
大きな曲がりはなく, 点はおおむね直線状に散らばっている.
極端に飛び離れた点も見当たらない.
直線を当てはめることに無理はなさそうだ.

\subsubsection*{Step 2 : 傾きと切片の計算}

% TODO: 具体的な計算を仮データで実施

$b$と$a$の公式に値を代入して計算する.
仮に$b = 2.3$, $a = 120$が得られたとする.

% [図挿入: 散布図に回帰直線を重ねたもの]

傾き$b = 2.3$は, 「広告費を1万円増やすと, 売上が平均的に2.3万円増える」と読める.
10万円増やせば, 平均的には23万円の売上増が見込める.
これが, 相関係数だけではわからなかった「量的な変化」の情報だ.

切片$a = 120$は, 式の上では「広告費がゼロのとき売上は120万円」となる.
しかし, データの広告費は20万円から80万円の範囲に収まっている.
広告費ゼロの月は実データに存在しない.
切片の120万円は, 直線の位置を決めるために計算上現れた数値であって, 「広告を出さなくても120万円売れる」という解釈にはそのまま使えない.

\subsubsection*{Step 3--4 : $R^2$と残差の確認}

$r = 0.87$だったから, $R^2 = 0.87^2 \approx 0.76$だ.
売上の散らばりの76\%が, 広告費との直線関係で説明できている.
残りの24\%は広告費だけでは説明しきれない.
季節, 天候, 競合の動きなど, 他の要因が影響しているのだろう.

残差プロットを描く.

% [図挿入: 残差プロット]

残差がゼロのまわりにランダムに散らばっていれば, 直線モデルはおおむね妥当だ.
もしU字型のパターンが見えたら, 広告費と売上の関係は直線ではなく曲線的かもしれない.
たとえば, ある広告費の水準を超えると効果が頭打ちになる, という可能性だ.

\subsubsection*{予測 : できることと, やってはいけないこと}

この直線を使えば, まだ見ていない月の売上を予測できる.
「来月の広告費を50万円にしたとき, 売上はおよそ$120 + 2.3 \times 50 = 235$万円」と見積もれる.
50万円はデータの範囲（20万円から80万円）に収まっているから, この予測にはそれなりの根拠がある.

では, 広告費200万円を代入して「売上580万円」と予測するのはどうか.
データの範囲をはるかに超えた値だ.
直線関係がそこまで続いている保証はない.
現実を考えても, 広告には効果の頭打ちがある.
予算を2倍にしたからといって売上が2倍になるとは限らない.
データの範囲外に直線を延ばして予測すること, いわゆる外挿は, 「直線が無限に続く」という仮定を暗に置いている.
その仮定が正しい保証はどこにもない.

前章の議論がここでも当てはまる.
「広告費と売上に強い直線関係がある」ことと, 「広告費が売上の原因である」ことは別の話だ.
$b = 2.3$は「広告費が1万円多い月は売上が平均2.3万円多い」という関係を記述しているだけであり, 「広告費を増やせば売上が増える」という因果を証明しているわけではない.
回帰直線は関係を記述し予測に使う道具であって, 因果を証明する道具ではない.

直線を引き, 数字を読み, 予測し, その限界を確認した.
何がわかって, 何がまだわからないのか.


\subsection{結果と次の一手 : 直線を引いた先に残る問い}

前章では, 2変数の関係の方向と強さを相関係数で測った.
本章では, その先の問い, 「$x$が変わったとき$y$はどのくらい変わるか」に答えた.

最小二乗法を使えば, 残差の二乗和が最小になる直線が一意に決まる.
誰が計算しても同じ$a$と$b$が得られるから, 目分量の問題は解消された.
傾き$b$は「$x$が1増えたときの$y$の平均的な変化量」を示し, $R^2$は直線がデータの散らばりをどれだけ説明しているかを0から1で教えてくれる.
残差プロットを確認すれば, 直線を信用してよいかどうかを自分の目で判断できる.
「関係がある」と言うだけでなく, 「どのくらい変わるか」を定量的に語れるようになったのだ.

ただし, ここまでの道具には前提がある.
データが独立に集められていること.
$x$と$y$の関係がおおむね直線的であること.
極端な外れ値がないこと.
前提が崩れていれば, 直線を引いて数字は出るが, その数字は当てにならない.

前提が崩れる場面は意外と多い.
標本の集め方に偏りがあれば, 回帰直線も偏る.
時間的に隣り合うデータが独立でなければ, 残差にパターンが現れ, $R^2$は説明力を過大に見積もるかもしれない.
因果の方向を取り違えていれば, 回帰の解釈そのものが間違う.
こうした問題は, 直線を引く技術の外側にある.
データの設計と解釈の問題だ.

次章では, 第2章から本章まで積み上げてきた手法の前提を整理し, その前提が崩れるとどうなるかを体系的に確認する.
バイアスの類型, 前提の確認方法, そしてデータを扱うときの倫理.
手法の「外側」を知ることが, 手法を正しく使うための最後の一歩になる.
